\chapter{Lezione 14 --- Formula di Ackermann e assegnamento dell'autostruttura}

\section{Impostazione del problema}

Consideriamo un sistema lineare tempo-invariante in forma di stato
\[
\begin{cases}
\dot x(t)=Ax(t)+Bu(t),\\[4pt]
y(t)=Cx(t),
\end{cases}
\]
con
\[
x(t)\in\mathbb{R}^n,\qquad
u(t)\in\mathbb{R}^m,\qquad
y(t)\in\mathbb{R}^\ell,
\]
e
\[
A\in\mathbb{R}^{n\times n},\qquad
B\in\mathbb{R}^{n\times m},\qquad
C\in\mathbb{R}^{\ell\times n}.
\]

Vogliamo assegnare la dinamica del sistema in ciclo chiuso mediante una
\textbf{retroazione completa dello stato}
\[
u(t)=-Kx(t),
\]
dove
\[
K\in\mathbb{R}^{m\times n}.
\]

Sostituendo nella dinamica di stato si ottiene
\[
\dot x(t)=Ax(t)+B(-Kx(t))=(A-BK)x(t).
\]

Quindi la matrice di stato del sistema retroazionato è
\[
A_{\text{cl}}=A-BK.
\]

Il problema consiste nello scegliere \(K\) in modo che \(A-BK\) abbia gli autovalori,
e possibilmente anche gli autovettori, desiderati.

\section{Richiamo: assegnamento degli autovalori}

Dalla lezione precedente sappiamo che:
\begin{itemize}
    \item nel caso SISO, se \((A,b)\) è completamente controllabile,
    si possono assegnare arbitrariamente tutti gli autovalori della matrice \(A-bk^T\);
    \item nel caso MIMO, se \((A,B)\) è completamente controllabile,
    si possono assegnare arbitrariamente gli autovalori della matrice \(A-BK\).
\end{itemize}

Nel caso MIMO, tuttavia, la soluzione non è unica, perché \(K\) ha più gradi di libertà.
Questo apre la possibilità di assegnare non solo gli autovalori, ma anche, entro certi limiti,
gli autovettori associati: cioè l'\textbf{autostruttura} del sistema in ciclo chiuso.

\section{Caso SISO: formula di Ackermann}

Consideriamo dapprima il caso SISO:
\[
\dot x=Ax+bu,\qquad y=c^T x,
\]
con
\[
b\in\mathbb{R}^n,\qquad c\in\mathbb{R}^n,
\]
e legge di retroazione
\[
u=-k^T x,
\qquad
k\in\mathbb{R}^n.
\]

La matrice chiusa è
\[
A-bk^T.
\]

\subsection{Polinomio caratteristico desiderato}

Supponiamo di voler assegnare gli autovalori desiderati
\[
\lambda_1,\lambda_2,\dots,\lambda_n.
\]

Il polinomio caratteristico desiderato è allora
\[
p_{\text{cl}}(\lambda)=
(\lambda-\lambda_1)(\lambda-\lambda_2)\cdots(\lambda-\lambda_n).
\]

Scrivendolo in forma espansa:
\[
p_{\text{cl}}(\lambda)=
\lambda^n+\alpha_{n-1}\lambda^{n-1}+\cdots+\alpha_1\lambda+\alpha_0.
\]

\subsection{Matrice di controllabilità}

Nel caso SISO la matrice di controllabilità è
\[
\mathcal{R}=
\begin{bmatrix}
b & Ab & A^2b & \cdots & A^{n-1}b
\end{bmatrix}.
\]

Se la coppia \((A,b)\) è completamente controllabile, allora \(\mathcal{R}\) è quadrata e invertibile:
\[
\mathcal{R}\in\mathbb{R}^{n\times n},
\qquad
\rank(\mathcal{R})=n.
\]

\subsection{Formula di Ackermann}

La \textbf{formula di Ackermann} permette di scrivere direttamente il guadagno di retroazione
senza passare esplicitamente per la forma canonica di controllo:
\[
k^T=
\begin{bmatrix}
0 & \cdots & 0 & 1
\end{bmatrix}
\mathcal{R}^{-1}\,p_{\text{cl}}(A),
\]
dove
\[
p_{\text{cl}}(A)=A^n+\alpha_{n-1}A^{n-1}+\cdots+\alpha_1A+\alpha_0 I.
\]

Questa formula è molto utile nel caso SISO, perché fornisce immediatamente il guadagno \(k\)
che assegna gli autovalori scelti.

\subsection*{Osservazione}

La formula di Ackermann è concettualmente equivalente al passaggio in forma canonica di controllo,
ma è più rapida dal punto di vista del calcolo simbolico o numerico.

\section{Caso MIMO: più gradi di libertà}

Passiamo ora al caso MIMO:
\[
\dot x=Ax+Bu,
\qquad
u=-Kx,
\qquad
K\in\mathbb{R}^{m\times n}.
\]

Anche qui la matrice chiusa è
\[
A-BK.
\]

Se la coppia \((A,B)\) è completamente controllabile,
allora è possibile assegnare arbitrariamente gli autovalori del sistema in ciclo chiuso.

Tuttavia, a differenza del caso SISO, qui:
\begin{itemize}
    \item la matrice \(K\) ha \(mn\) coefficienti;
    \item il numero di vincoli imposti dal solo polinomio caratteristico è \(n\).
\end{itemize}

Quindi, in generale, rimangono dei gradi di libertà.
Questi possono essere sfruttati per cercare di assegnare anche gli autovettori,
cioè l'autostruttura del sistema chiuso.

\section{Assegnamento dell'autostruttura}

Supponiamo di voler assegnare a \(A-BK\) gli autovalori
\[
\lambda_1,\lambda_2,\dots,\lambda_n
\]
e, per quanto possibile, anche gli autovettori
\[
v_1,v_2,\dots,v_n.
\]

Se \(\lambda_i\) è autovalore di \(A-BK\) e \(v_i\) è il corrispondente autovettore,
allora deve valere
\[
(A-BK)v_i=\lambda_i v_i.
\]

Portando tutto a sinistra:
\[
(\lambda_i I-A+B K)v_i=0.
\]

Scrivendo
\[
q_i:=Kv_i\in\mathbb{R}^m,
\]
la relazione precedente diventa
\[
(\lambda_i I-A)v_i+Bq_i=0.
\]

Questa equazione si può riscrivere in forma matriciale come
\[
\begin{bmatrix}
\lambda_i I-A & B
\end{bmatrix}
\begin{bmatrix}
v_i\\
q_i
\end{bmatrix}
=0.
\]

Definiamo quindi
\[
T_i:=
\begin{bmatrix}
\lambda_i I-A & B
\end{bmatrix}
\in\mathbb{R}^{n\times (n+m)}.
\]

Dobbiamo allora cercare vettori
\[
p_i=
\begin{bmatrix}
v_i\\
q_i
\end{bmatrix}
\in \ker(T_i).
\]

\section{Dimensione del nucleo di \(T_i\)}

Poiché \((A,B)\) è completamente controllabile, per il criterio di Hautus si ha
\[
\rank\!\begin{bmatrix}
\lambda I-A & B
\end{bmatrix}=n
\qquad
\forall \lambda\in\mathbb{C}.
\]

Quindi, per ogni \(\lambda_i\),
\[
\rank(T_i)=n.
\]

Essendo \(T_i\in\mathbb{R}^{n\times (n+m)}\), segue
\[
\dim\ker(T_i)=(n+m)-n=m.
\]

Dunque, per ciascun autovalore scelto \(\lambda_i\),
esistono \(m\) vettori linearmente indipendenti nel nucleo di \(T_i\).

\section{Costruzione dei vettori \(p_i\)}

Sia
\[
\{m_{i1},m_{i2},\dots,m_{im}\}
\]
una base di \(\ker(T_i)\).

Allora ogni vettore del nucleo si può scrivere come combinazione lineare
\[
p_i=\sum_{j=1}^{m}\alpha_{ij}\,m_{ij}.
\]

Poiché vogliamo che le prime \(n\) componenti di \(p_i\) coincidano con l'autovettore desiderato \(v_i\),
scriviamo
\[
p_i=
\begin{bmatrix}
v_i\\
q_i
\end{bmatrix}.
\]

La scelta dei coefficienti \(\alpha_{ij}\) viene fatta proprio in modo che
la parte alta del vettore \(p_i\) sia il vettore \(v_i\) desiderato.
La parte bassa fornisce automaticamente il vettore
\[
q_i=Kv_i.
\]

Ripetendo questa procedura per tutti gli autovalori \(\lambda_i\),
si ottengono i vettori
\[
p_1,p_2,\dots,p_n.
\]

\section{Costruzione di \(K\)}

A questo punto costruiamo le matrici
\[
V=
\begin{bmatrix}
v_1 & v_2 & \cdots & v_n
\end{bmatrix},
\qquad
Q=
\begin{bmatrix}
q_1 & q_2 & \cdots & q_n
\end{bmatrix}.
\]

Poiché per ciascun \(i\) vale
\[
q_i=Kv_i,
\]
mettendo insieme tutte le colonne si ottiene
\[
KV=Q.
\]

Se \(V\) è invertibile, allora
\[
K=QV^{-1}.
\]

Questa è la formula finale che permette di ricavare la matrice di guadagno
a partire dagli autovettori desiderati e dai corrispondenti vettori \(q_i\).

\section{Algoritmo operativo}

Riassumendo, l'algoritmo per l'assegnamento dell'autostruttura nel caso MIMO è:

\begin{enumerate}
    \item fissare gli autovalori desiderati
    \[
    \lambda_1,\lambda_2,\dots,\lambda_n;
    \]

    \item per ogni \(i=1,\dots,n\), costruire la matrice
    \[
    T_i=
    \begin{bmatrix}
    \lambda_i I-A & B
    \end{bmatrix};
    \]

    \item calcolare una base del nucleo \(\ker(T_i)\);

    \item scegliere una combinazione lineare dei vettori della base in modo che
    la parte alta del vettore risultante sia l'autovettore desiderato \(v_i\);

    \item leggere la parte bassa come \(q_i\), cioè
    \[
    p_i=
    \begin{bmatrix}
    v_i\\
    q_i
    \end{bmatrix};
    \]

    \item costruire le matrici
    \[
    V=
    \begin{bmatrix}
    v_1 & \cdots & v_n
    \end{bmatrix},
    \qquad
    Q=
    \begin{bmatrix}
    q_1 & \cdots & q_n
    \end{bmatrix};
    \]

    \item se \(V\) è invertibile, porre
    \[
    K=QV^{-1}.
    \]
\end{enumerate}

\section{Esempio}

Consideriamo il sistema
\[
\dot x=
\begin{bmatrix}
-1 & 1\\
0 & -1
\end{bmatrix}x
+
\begin{bmatrix}
1 & 0\\
0 & 1
\end{bmatrix}u.
\]

Qui:
\[
n=2,\qquad m=2.
\]

Vogliamo assegnare in ciclo chiuso gli autovalori
\[
\lambda_1=-5,\qquad \lambda_2=-10.
\]

\subsection{Primo autovalore}

Per \(\lambda_1=-5\):
\[
T_1=
\begin{bmatrix}
-4 & -1 & 1 & 0\\
0 & -4 & 0 & 1
\end{bmatrix}.
\]

Una base del nucleo è, ad esempio,
\[
m_{11}=
\begin{bmatrix}
1\\0\\4\\0
\end{bmatrix},
\qquad
m_{12}=
\begin{bmatrix}
0\\1\\1\\4
\end{bmatrix}.
\]

Scegliamo come autovettore desiderato
\[
v_1=
\begin{bmatrix}
1\\0
\end{bmatrix}.
\]

Allora basta prendere
\[
p_1=m_{11}=
\begin{bmatrix}
1\\0\\4\\0
\end{bmatrix},
\]
da cui
\[
q_1=
\begin{bmatrix}
4\\0
\end{bmatrix}.
\]

\subsection{Secondo autovalore}

Per \(\lambda_2=-10\):
\[
T_2=
\begin{bmatrix}
-9 & -1 & 1 & 0\\
0 & -9 & 0 & 1
\end{bmatrix}.
\]

Una base del nucleo è
\[
m_{21}=
\begin{bmatrix}
1\\0\\9\\0
\end{bmatrix},
\qquad
m_{22}=
\begin{bmatrix}
0\\1\\1\\9
\end{bmatrix}.
\]

Scegliamo come autovettore desiderato
\[
v_2=
\begin{bmatrix}
0\\1
\end{bmatrix}.
\]

Allora possiamo prendere
\[
p_2=m_{22}=
\begin{bmatrix}
0\\1\\1\\9
\end{bmatrix},
\]
da cui
\[
q_2=
\begin{bmatrix}
1\\9
\end{bmatrix}.
\]

\subsection{Costruzione finale di \(K\)}

Si ottiene quindi
\[
V=
\begin{bmatrix}
1 & 0\\
0 & 1
\end{bmatrix},
\qquad
Q=
\begin{bmatrix}
4 & 1\\
0 & 9
\end{bmatrix}.
\]

Poiché \(V=I\), segue immediatamente
\[
K=QV^{-1}=Q=
\begin{bmatrix}
4 & 1\\
0 & 9
\end{bmatrix}.
\]

Verifica:
\[
A-BK=
\begin{bmatrix}
-1 & 1\\
0 & -1
\end{bmatrix}
-
\begin{bmatrix}
4 & 1\\
0 & 9
\end{bmatrix}
=
\begin{bmatrix}
-5 & 0\\
0 & -10
\end{bmatrix}.
\]

Quindi gli autovalori assegnati sono effettivamente
\[
-5,\qquad -10,
\]
e in questo caso si è assegnata completamente anche l'autostruttura.

\section{Osservazioni numeriche e robustezza}

Nell'esempio precedente l'inversione di \(V\) è banale.
In generale, però, la costruzione
\[
K=QV^{-1}
\]
può essere numericamente delicata.

Se \(V\) è invertibile ma mal condizionata, cioè con determinante molto piccolo o colonne quasi dipendenti,
l'inversione numerica può introdurre errori rilevanti.
In tal caso il guadagno calcolato potrebbe non collocare gli autovalori esattamente dove previsto.

Per migliorare la robustezza numerica conviene, quando possibile, scegliere autovettori
il più possibile ortogonali o comunque ben condizionati.

\section{Quando si può assegnare tutta l'autostruttura?}

Nel caso dell'esempio avevamo
\[
n=m.
\]

Quindi il numero di ingressi indipendenti era uguale al numero di stati,
e questo ha permesso di assegnare completamente l'autostruttura:
non solo gli autovalori, ma anche tutti gli autovettori scelti.

In generale:
\begin{itemize}
    \item gli autovalori si possono assegnare se \((A,B)\) è completamente controllabile;
    \item l'assegnamento completo degli autovettori dipende dal numero di ingressi e dalla struttura del sistema;
    \item in MIMO ci sono più gradi di libertà, ma non sempre è possibile assegnare arbitrariamente \emph{ogni} autovettore.
\end{itemize}

\section{Caso tempo discreto}

Negli appunti si osserva che nei sistemi a tempo discreto tutto funziona in modo del tutto analogo.

Consideriamo
\[
x(k+1)=Ax(k)+Bu(k),
\]
con retroazione di stato
\[
u(k)=-Kx(k).
\]

Il sistema chiuso diventa
\[
x(k+1)=(A-BK)x(k).
\]

La teoria, le formule e la procedura sono le stesse del caso continuo,
con l'unica differenza concettuale che nel tempo discreto l'ipotesi corretta è che
la coppia \((A,B)\) sia \textbf{completamente raggiungibile}.

Nel caso tempo discreto:
\begin{itemize}
    \item si assegnano autovalori desiderati nel piano \(z\);
    \item per la stabilità, tali autovalori devono stare dentro il cerchio unitario;
    \item la formula di Ackermann e l'algoritmo MIMO si applicano nello stesso modo,
    sostituendo al polinomio in \(\lambda\) il polinomio desiderato in \(z\).
\end{itemize}

\section{Sintesi della lezione}

I punti principali della lezione sono:
\begin{enumerate}
    \item nel caso SISO, il guadagno di retroazione può essere calcolato direttamente con la formula di Ackermann:
    \[
    k^T=
    \begin{bmatrix}
    0 & \cdots & 0 & 1
    \end{bmatrix}
    \mathcal{R}^{-1}p_{\text{cl}}(A);
    \]

    \item nel caso MIMO, la soluzione non è unica perché \(K\) ha più gradi di libertà;

    \item questi gradi di libertà possono essere sfruttati per assegnare non solo gli autovalori,
    ma anche gli autovettori, cioè l'autostruttura;

    \item per ogni autovalore desiderato \(\lambda_i\), si costruisce
    \[
    T_i=
    \begin{bmatrix}
    \lambda_i I-A & B
    \end{bmatrix};
    \]

    \item il vettore
    \[
    p_i=
    \begin{bmatrix}
    v_i\\
    q_i
    \end{bmatrix}
    \]
    deve appartenere al nucleo di \(T_i\);

    \item ripetendo la costruzione per tutti gli autovalori si ottengono
    \[
    KV=Q,
    \qquad
    K=QV^{-1};
    \]

    \item in presenza di matrici mal condizionate, l'inversione di \(V\) può creare problemi numerici;

    \item nel caso tempo discreto la teoria è la stessa del continuo, purché la coppia \((A,B)\) sia completamente raggiungibile.
\end{enumerate}